\section{游戏环境与公共实验基础}
\label{sec:shared}
\subsection{游戏规则与物理过程}
项目在 Gymnasium 接口下实现固定球速 Breakout 环境。画面分辨率为 $210\times160$，包含小球、挡板、砖块、边界以及分数和剩余生命显示。每局有 5 条生命，每面砖墙有 $6\times18=108$ 块砖，共设置两面砖墙。各行砖块从上至下分别计 7、7、4、4、1、1 分，因此一面墙为 432 分，整局最高分为 864 分。

物理系统以 60 Hz 更新，小球每帧移动 2 像素，即球速固定为 120 像素/秒。挡板宽度固定为 16 像素，单个物理帧向左或向右移动 3 像素。碰撞后小球改变运动方向但保持速度模长。将球速和挡板宽度固定，可以减少随游戏进程变化的额外难度因素，使模型在一致的控制条件下学习。

小球与挡板接触的位置决定反弹方向。挡板横向划分为五个区域，越靠近边缘，反弹方向越偏向相应侧；中心区域保留小角度水平分量，避免长时间完全垂直往返。该规则不随累计接球次数改变。砖块碰撞使用运动路径与扩展矩形的相交检测，在一个物理步内寻找最早接触点，处理反射并继续推进剩余时间，以避免仅检查离散终点造成漏碰撞。

发球位置和初始水平方向由随机种子控制。第一面砖墙清空后重建第二面墙并重新发球；生命耗尽或第二面砖墙清空时游戏真正结束。该环境是按任务需要重建的游戏，其几何和反弹规则由项目代码定义，实验分数对应这一固定规则环境。

\reportfigure[width=.72\linewidth,height=9cm,keepaspectratio]{images/game_live.png}{游戏实时展示画面}{网页运行截图，展示游戏区域和演示控制。}{DQN 智能体在共用网页展示接口中的实时运行画面，显示游戏区域、得分、剩余生命与演示控制。}{fig:game}

\subsection{状态、动作与奖励定义}
物理环境的完整状态可以表示为小球位置与速度、挡板位置、砖块存活标记、剩余生命、当前墙面以及发球状态等变量的集合。这些变量决定下一步物理变化。智能体实际获得的是渲染图像，模型不直接读取球坐标、速度或预先计算的落点。

动作空间为四个离散动作，如表\ref{tab:actions}所示。自动发球包装器在等待发球时将当前动作替换为发球动作，并在正常的物理步中执行，不额外插入隐藏帧。因此，学习任务主要集中在接球与击砖控制。

\begin{table}[!htbp]
 \centering
 \caption{环境动作空间与控制含义}\label{tab:actions}
 \begin{tabularx}{.86\linewidth}{l Y}
 \toprule
 动作名称 & 物理控制含义\\
 \midrule
 NOOP & 挡板保持当前位置\\
 FIRE & 等待发球时发射小球\\
 RIGHT & 挡板向右移动，受右边界约束\\
 LEFT & 挡板向左移动，受左边界约束\\
 \bottomrule
 \end{tabularx}
\end{table}

奖励直接采用当前动作执行期间击碎砖块获得的原生分数：
\begin{equation}
 r_t=\sum_{b\in\mathcal B_t}v(b),
\end{equation}
其中 $\mathcal B_t$ 为该决策步内被击碎的砖块集合，$v(b)$ 为砖块所在行对应的分值。DQN 和最终 PPO 对照均采用原始奖励，不设置接球奖励、存活奖励或失球惩罚，也不裁剪奖励。辅助奖励的单独消融见第\ref{sec:ppo}节。所有评估只累计原生得分。失球虽无直接负奖励，但会减少剩余生命和未来得分机会。

\subsection{图像预处理与时序信息}
单张画面可以反映小球的位置，却无法可靠区分运动方向。项目将连续四张处理后的画面堆叠，作为对运动状态的近似表示。共享环境工厂按以下顺序处理图像。

首先，将一个动作重复执行 4 个物理帧，并累加期间的奖励。随后，对实际产生的最后两帧逐像素取最大值，得到池化画面。若提前终止或截断，则使用真正发生的末尾画面，避免把旧帧当成终止观测。之后，使用灰度转换和缩放得到 $84\times84$ 图像，再按时间顺序堆叠 4 帧，并将通道维移动到最前端。

因此，单个状态表示为
\begin{equation}
 s_t=[o_{t-3},o_{t-2},o_{t-1},o_t]
 \in\{0,\ldots,255\}^{4\times84\times84}.
\end{equation}
单环境批输入的形状为 $(1,4,84,84)$，原始存储类型为 \texttt{uint8}。网络内部将像素转换为浮点数并除以 255。堆帧允许卷积网络利用位置变化判断运动趋势，但它仍是从图像得到的近似状态表示。

\reportfigure{images/observation_pipeline.png}{图像预处理示意}{原始 RGB 画面、最大池化结果与四帧灰度输入。}{共享环境工厂实际采集的一组连续观测。四张灰度图按时间顺序组成 $4\times84\times84$ 输入，最新一帧已逐像素核对实际网络观测。}{fig:preprocess}

\subsection{回合终止、时间截断与环境一致性}
项目以一整局游戏作为一个回合，损失一条生命不会立即终止回合。只要仍有剩余生命，环境就准备并自动发射下一球。这样可以让价值估计覆盖一局中的多次失球和后续得分机会。

所有模型使用统一的单局时间上限。达到上限属于时间截断，生命耗尽或通关属于真正终止，两者在学习目标中需要区分。真正终止时未来价值为零；时间截断则仍需要基于截断时的最终观测进行适当自举。环境包装器保留相应标记和终止观测，由训练库的回放与目标计算流程处理。

训练、独立评估、实时展示和视频录制共用同一套环境与预处理代码。模型保存环境版本和关键参数，加载时检查球速、挡板宽度、图像尺寸、堆帧数量和奖励规则等设置。该检查能够阻止外形相同但物理规则不同的历史模型被误用于当前任务。

\subsection{强化学习任务与长期回报}
设 $s_t$ 为时刻 $t$ 提供给智能体的状态表示，$a_t$ 为动作，$r_t$ 为执行该动作后获得的奖励，$T$ 为本局结束时刻。折扣累计回报定义为
\begin{equation}
 G_t=\sum_{k=0}^{T-t-1}\gamma^k r_{t+k},\qquad
 J(\pi)=\mathbb{E}_{\pi}[G_0],
\end{equation}
其中 $\pi$ 为策略，$\gamma\in[0,1)$ 为折扣因子。DQN 与初始 PPO 取 $\gamma=0.99$，后续 PPO 实验取 0.995，以调整对远期击砖收益的权重。训练目标是提高期望折扣回报；实验评估则报告整局未折扣的游戏总分，两者含义不同。

\subsection{共用系统结构与实时展示}
系统使用 Python、Gymnasium 和 Stable-Baselines3 实现 DQN 与 PPO。游戏、预处理、参数管理、模型加载、评估和实时推理接口共用。算法各自的采样、更新及续训机制分别在第\ref{sec:dqn}节和第\ref{sec:ppo}节说明。

环境层接收动作，推进物理过程并返回图像、奖励和结束标记；管理层定期评估并保存检查点；展示层用相同观测接口确定性推理。各层共享任务定义，训练和评估使用独立环境实例。

\begin{table}[!htbp]
 \centering\small
 \caption{共用程序模块及其职责}\label{tab:modules}
 \begin{tabularx}{\linewidth}{P{4.5cm}Y}
 \toprule
 模块 & 主要实现内容\\
 \midrule
 \texttt{fixed\_breakout.py} & 游戏物理、碰撞、计分、渲染、自动发球与动作重复\\
 \texttt{env.py} & 环境工厂、版本检查、灰度缩放、帧堆叠与通道转换\\
 \texttt{model\_io.py} & 识别 DQN/PPO 变体并检查游戏契约\\
 \texttt{callbacks.py} & 开发评估、最佳模型选择、早停与训练状态记录\\
 \texttt{evaluate.py} & 逐局种子评估、结果导出及随机策略比较入口\\
 \texttt{dashboard.py}、\texttt{game\_viewer.py} & 网页展示、模型切换、暂停及播放控制\\
 \bottomrule
 \end{tabularx}
\end{table}

展示采用确定性动作：DQN 选择最大 Q 值，PPO 选择最大动作概率。网页提供模型选择、随机策略、暂停、倍速和全屏，显示实际物理帧，1 倍以 60 帧/秒为目标。定量评价仍依赖多局统计。

\subsection{评价指标与数据使用原则}
每局游戏总分定义为 $R_i=\sum_{t=0}^{T_i-1}r_t$。对 $N$ 局游戏，平均分与总体形式的标准差分别为
\begin{equation}
 \overline R=\frac{1}{N}\sum_{i=1}^{N}R_i,\qquad
 \sigma_R=\sqrt{\frac{1}{N}\sum_{i=1}^{N}(R_i-\overline R)^2}.
\end{equation}
平均分衡量总体得分水平，中位数反映典型游戏表现，标准差和低分位数用于分析波动。仅报告最高分容易放大少数成功轨迹的影响；仅报告损失则不能直接反映控制效果。

模型选择使用开发数据，冻结候选后再汇总统一测试结果。相同开局种子上的逐局分数用于配对比较；区间按完整种子对进行 bootstrap 重采样。此类区间反映固定模型在回合样本上的差异，不覆盖独立重复训练的波动。不同算法的损失目标和尺度不同，损失只用于各自的训练诊断。

\begin{table}[!htbp]
 \centering\footnotesize
 \caption{报告中的数据用途与评估协议}\label{tab:data-protocols}
 \begin{tabularx}{\linewidth}{P{3.1cm}P{3.8cm}Y}
 \toprule
 数据用途 & 采样与隔离方式 & 报告位置与解释\\
 \midrule
 DQN 训练开发评估 & 固定随机性设置，连续生成游戏序列 & 第\ref{sec:dqn}节；用于各阶段选模\\
 DQN 开发复测 & 逐局重置，各候选共享开局集合 & 第\ref{sec:dqn}节；30 局候选诊断\\
 PPO 起点选择与开发 & 起点验证与训练开发使用不同集合 & 第\ref{sec:ppo}节；分别用于起点选择和训练选模\\
 历轮 PPO 实验评估 & 各轮使用各自记录的种子集合 & 第\ref{sec:ppo}节；不拼接为同测试集提升曲线\\
 冻结模型统一测试 & 使用留出的共同开局，逐局重置 & 第\ref{sec:comparison}节；50 局方法比较与最终 GAE 消融\\
 \bottomrule
 \end{tabularx}
\end{table}


\FloatBarrier
